\documentclass[11pt,reqno]{amsart}
\usepackage[T1]{fontenc}
\usepackage{lmodern}
\usepackage[a4paper,margin=28mm]{geometry}
\usepackage[expansion=false]{microtype}
\usepackage{booktabs}
\usepackage{float}
\usepackage[hidelinks]{hyperref}
\hypersetup{pdftitle={Four-Cycle-Free Subgraphs of the Seven-Dimensional Hypercube}}

\numberwithin{equation}{section}
\newtheorem{theorem}{Theorem}[section]
\newtheorem{lemma}[theorem]{Lemma}
\newtheorem{proposition}[theorem]{Proposition}
\newcommand{\ex}{\operatorname{ex}}
\newcommand{\one}{\mathbf{1}}

\title[Four-cycle-free subgraphs of the seven-cube]
  {Four-Cycle-Free Subgraphs of the Seven-Dimensional Hypercube}
% Add author names and affiliations here.
\thanks{Working draft.}

\begin{document}

\begin{abstract}
Let \(\ex(Q_n,C_4)\) denote the largest number of edges in a
\(C_4\)-free subgraph of the \(n\)-dimensional hypercube. We prove that
\(\ex(Q_7,C_4)=304\). In a hypothetical \(305\)-edge subgraph, counting
three-edge stars on the two parity classes forces every degree to be
four or five. A square-counting identity and a relation between opposite
facets then reduce the proof to five possible exceptions to a local
inequality. The three six-dimensional obstructions used in the proof
are certified by \(21\) independently checked propositional refutations.
\end{abstract}

\maketitle

\section{Introduction}\label{sec:intro}

The hypercube \(Q_n\) has vertex set \(\{0,1\}^n\), with two vertices
adjacent when they differ in exactly one coordinate. Following
Erd\H{o}s~\cite{erdos}, we consider the extremal function
\[
  f(n)=\ex(Q_n,C_4)
       =\max\{e(G):G\subseteq Q_n,\ G\text{ is }C_4\text{-free}\}.
\]
Brass, Harborth, and Nienborg~\cite{bhn} studied finite-dimensional
bounds and constructions. Harborth and Nienborg~\cite{hn} determined
\(f(6)=132\).

An odd-square edge set meets every coordinate square in one or three
edges and is therefore \(C_4\)-free. The fully frustrated hypercube
construction of Derrida, Pomeau, Toulouse, and Vannimenus~\cite{dptv}
attains \(304\) edges in dimension seven. Minamoto~\cite{minamoto} gives
a machine-verifiable edge list and proves that \(304\) is optimal
among odd-square subgraphs of \(Q_7\). Minamoto's Open Problem~1 asks
whether \(304\) is also optimal in the unrestricted class. We answer
this question affirmatively.

\begin{theorem}\label{thm:q7}
The maximum number of edges in a \(C_4\)-free subgraph of \(Q_7\) is
\(304\).
\end{theorem}

For the upper bound, we suppose a \(305\)-edge example exists and
restrict it to coordinate faces. The central inequality relates the
edge count of a six-dimensional facet to the difference between the
number of empty squares and the number of squares containing exactly
two opposite edges. Counting across opposite facets leaves five
possible profiles for a violation of this inequality. The required
local exclusions use \(21\) propositional refutations.

Section~\ref{sec:prelims} gives the basic counting identities.
Section~\ref{sec:q7} proves Theorem~\ref{thm:q7}.
Appendix~\ref{app:certificates} specifies the finite encodings, the
symmetry reduction, and the proof checks used in that argument.

\section{Coordinate faces and counting identities}\label{sec:prelims}

We treat subgraphs as spanning by allowing isolated vertices. Let
\(P_0\) and \(P_1\) be the two classes of vertices according to the parity
of their Hamming weight. A \(k\)-dimensional coordinate face is obtained
by choosing \(k\) free coordinates and fixing the remaining coordinates.
We use \(G[F]\) for the restriction of a graph \(G\) to a face \(F\).
When the ambient face is understood, we use the same letter for its
restriction. An \((n-1)\)-dimensional coordinate face of \(Q_n\) is a
\emph{facet}. Write \(d_G(v)\) for degree and \(e(G)\) for edge count.

\begin{lemma}[Face multiplicities]\label{lem:multiplicities}
There are \(\binom nk 2^{n-k}\) coordinate \(k\)-faces of \(Q_n\), and
each coordinate \(r\)-face lies in \(\binom{n-r}{k-r}\) of them for
\(r\le k\le n\).
\end{lemma}

\begin{proof}
Choose the \(k\) free coordinates and assign values to the remaining
\(n-k\). To extend an \(r\)-face to a \(k\)-face, choose \(k-r\) of
its fixed coordinates to free.
\end{proof}

A \emph{full three-star} consists of three retained edges incident
with a common vertex, its root. The three directions determine a
unique coordinate \(3\)-face containing the star. For a parity class
\(P\), write
\[
  S_P(G)=\sum_{v\in P}\binom{d_G(v)}{3}.
\]

\begin{lemma}[Star capacity]\label{lem:star}
In a \(C_4\)-free graph, each coordinate \(3\)-face contains at most
one full three-star rooted on a specified parity side. Consequently
\[
  S_P(G)\le \binom{n}{3}2^{n-3}.
\]
The difference counts the \(3\)-faces \emph{missing a \(P\)-star}.
\end{lemma}

\begin{proof}
Two distinct same-parity vertices of a \(3\)-face differ in exactly
two coordinates. If both were roots of full stars in the face, the four
edges in those two directions incident with the two roots would form a
\(C_4\). Summing over the \(3\)-faces proves the bound and the counting
interpretation.
\end{proof}

Write \(a_0(G)\) for the number of empty coordinate squares and
\(a_{\mathrm{opp}}(G)\) for the number with exactly two retained edges
that are opposite. Define \(D(G)=a_0(G)-a_{\mathrm{opp}}(G)\).

\begin{lemma}[Square identity]\label{lem:square}
If \(G\subseteq Q_n\) is \(C_4\)-free, then
\begin{equation}\label{eq:square}
  \sum_{v}d_G(v)^2
    =2n\,e(G)-2\binom n2 2^{n-2}+2D(G).
\end{equation}
\end{lemma}

\begin{proof}
Every adjacent pair of retained edges belongs to a unique square. In a
square with \(k\) retained edges, the number of adjacent pairs is
\(k-1+\one_{\{\text{empty}\}}-\one_{\{\text{two opposite edges}\}}\),
as is seen by considering \(k=0,1,2,3\). Each retained edge lies in
\(n-1\) squares. Summing over squares gives
\(\sum_v\binom{d_G(v)}2=(n-1)e(G)-\binom n2 2^{n-2}+D(G)\).
Now use \(d^2=2\binom d2+d\) and \(\sum_v d_G(v)=2e(G)\).
\end{proof}

\section{The exact extremal number in dimension seven}\label{sec:q7}

\subsection{Local obstructions in dimension six}

\begin{lemma}[A single degree-two vertex]\label{lem:root}
There is no \(C_4\)-free subgraph of \(Q_6\) in which one parity side
has exactly one vertex of degree two, all its other vertices have
degree in \(\{3,4,5\}\), and its star count is \(160\).
\end{lemma}

\begin{proof}
Translate the degree-two vertex to the origin and permute coordinates
so that its two retained directions are \(0\) and \(1\). Use a Boolean
variable for each of the \(192\) edges and an indicator for each of
the \(640\) candidate full stars rooted on the even side. For each
square add a clause forbidding its four edges from all being retained.
Set each star indicator equal to the conjunction of its three edges.
The star-count hypothesis and Lemma~\ref{lem:star} require exactly
one even-rooted star in each of the \(160\) three-dimensional faces.
Fix the six incident edge variables at the origin and impose the
stated degree bounds at the remaining \(31\) even vertices.

A graph in the statement would satisfy this formula, with the star
indicators set by its edges. The resulting CNF has \(832\) variables
and \(4422\) clauses. The proof check in
Appendix~\ref{app:replay} verifies a refutation of this formula.
\end{proof}

\begin{lemma}[Two star-count bounds]\label{lem:rigid}
Let \(H\subseteq Q_6\) be \(C_4\)-free, with all degrees in
\(\{3,4,5\}\).
\begin{enumerate}
\item[\(\mathrm{R0}\).] If the two parity star counts are \(158,158\),
  then \(e(H)\ge132\).
\item[\(\mathrm{R1}\).] If the two parity star counts are \(158,155\)
  in either order, then \(e(H)\ge130\).
\end{enumerate}
\end{lemma}

\begin{proof}
On a parity side with \(b\) vertices of degree three,
\begin{equation}\label{eq:local-stars}
  S=6e(H)-640+3b.
\end{equation}
Indeed, \(\binom d3=6d-20+3\one_{\{d=3\}}\) for \(d\in\{3,4,5\}\),
and the side has \(32\) vertices. Thus on a side with \(S=158\) one
has \(b=266-2e(H)\). A counterexample to \(\mathrm{R0}\) has \(b\ge4\)
on that side, and a counterexample to \(\mathrm{R1}\) has \(b\ge8\).

Each \(Q_6\) has \(160\) coordinate \(3\)-faces, so this side misses
a star in exactly two of them. Normalize the first such face to the
face with free coordinates \(0,1,2\) and all other coordinates zero.
This also allows us to use even parity for the side of interest.
The stabilizer of this face preserving even parity has ten orbits on
the remaining faces, derived in Appendix~\ref{app:orbits}. For each
orbit, require the stated two missing even stars, the appropriate
number of missing odd stars (two or five), the degree bounds, and the
lower bound on \(b\). Each counterexample therefore extends to one
of the ten Boolean inputs for the corresponding assertion.
Appendices~\ref{app:encoding} and~\ref{app:replay} give the encoding
and the refutations of all twenty inputs.
\end{proof}

\subsection{Degrees in a hypothetical 305-edge graph}

If a counterexample to the upper bound contains more than \(305\)
edges, delete edges until exactly \(305\) remain. We therefore fix a
\(C_4\)-free graph \(G\subseteq Q_7\) with \(e(G)=305\). Throughout this
subsection, counts \(a,b,c\) refer to one chosen parity side \(P\).

\begin{lemma}[Initial degree restrictions]\label{lem:initial}
On each parity side of \(G\), the degrees belong to \(\{3,4,5\}\)
and at most two vertices have degree three.
\end{lemma}

\begin{proof}
Let \(\omega(d)=\binom d3-6d+20\). Its values at \(d=0,\ldots,7\) are
\(20,14,8,3,0,0,4,13\), respectively. Since each side has \(64\)
vertices with degree sum \(305\), Lemma~\ref{lem:star} gives
\[
  \sum_{v\in P}\omega(d_G(v))=S_P(G)-550\le10 .
\]
In particular, degrees \(0,1,7\) are excluded. Let \(a,b,c\) count
vertices of degrees \(2,3,6\) in \(P\), respectively. The penalty
bound and the degree sum yield
\begin{equation}\label{eq:penalty}
  8a+3b+4c\le10,\qquad n_4(P)=15-3a-2b+c.
\end{equation}

For each \(6\)-facet \(H\), put \(\delta_H=132-e(H)\) and
let \(A_H,B_H,C_H\) count degrees \(1,2,6\) in \(P\cap V(H)\)
after restriction. Write \(g(d)=(d-4)(d-5)/2\). For the possible
local degrees \(1\le d\le6\), one has
\[
  \binom d3=3g(d)+6d-20
      -4\one_{\{d=1\}}-\one_{\{d=2\}}+\one_{\{d=6\}} .
\]
The local degree sum on one side is \(e(H)\), and its star count is
at most \(160\). Rearranging and rounding down gives
\begin{equation}\label{eq:floor}
  \sum_{v\in P\cap V(H)}g(d_H(v))
  \le 2+2\delta_H+
      \left\lfloor\frac{2+4A_H+B_H-C_H}{3}\right\rfloor .
\end{equation}

A vertex of degree \(d\) occurs in seven facets, with local degree
\(d-1\) in \(d\) of them and \(d\) in the remaining \(7-d\).
Its contribution to the sum of the left-hand sides is
\(d\,g(d-1)+(7-d)g(d)\), taking values \(27,13,4,0,1\)
for \(d=2,3,4,5,6\). By \eqref{eq:penalty} the total is
\(60+15a+5b+5c\). Each edge belongs to six facets, so
\(\sum_H\delta_H=14\cdot132-6\cdot305=18\).

If \(a=0\), then \(A_H=0\) and each floor in \eqref{eq:floor}
is at most \(B_H\). Each global degree-three vertex contributes to
\(B_H\) in three facets, giving \(\sum_H B_H=3b\).
If \(a=1\), then \eqref{eq:penalty} gives \(b=c=0\). The sole
degree-two vertex contributes \(A_H=1\) in two facets and
\(B_H=1\) in five others, and the sum of the floors is \(9\).
These cases exhaust \(a\) by \eqref{eq:penalty}. Summing
\eqref{eq:floor} gives
\[
  60+15a+5b+5c\le 64+9a+3b,
  \qquad 6a+2b+5c\le4 .
\]
Thus \(a=c=0\) and \(b\le2\) on the chosen side. The argument
applies to either parity side.
\end{proof}

\begin{lemma}[The degree sequence]\label{lem:degrees}
Each parity side of \(G\) has \(15\) vertices of degree four and
\(49\) of degree five.
\end{lemma}

\begin{proof}
Fix one parity side \(P\). By Lemma~\ref{lem:initial} the number \(b\)
of degree-three vertices on this side belongs to \(\{0,1,2\}\).

If \(b=2\), equality holds throughout the summed bounds in the proof
of Lemma~\ref{lem:initial}. In particular,
\(\lfloor(2+B_H)/3\rfloor=B_H\) for every facet, so \(B_H\in\{0,1\}\).
Since \(\sum_H B_H=3b=6\), six facets have \(B_H=1\). In each such
facet equality in \eqref{eq:floor} implies a star count of \(160\).
Indeed, \(2+B_H=3\) is divisible by three, so the rounding step
loses nothing there. The other \(31\) vertices on that side have
local degree in \(\{3,4,5\}\). This contradicts
Lemma~\ref{lem:root}.

If \(b=1\), the unique degree-three vertex of \(P\) belongs to seven
facets. It drops to degree two in three and retains degree three in
four. The seven facets omitting the vertex and the four in which it
retains its degree all have their local degrees on \(P\) in
\(\{3,4,5\}\). Since
\(\binom33\equiv\binom43\equiv\binom53\equiv1\pmod3\),
these eleven facets have star counts congruent to \(32\pmod3\),
hence at most \(158\).
In the three other facets, the star count is congruent to \(31\pmod3\),
and Lemma~\ref{lem:root} excludes \(160\), so it is at most \(157\).
But \(S_P(G)=6\cdot305-20\cdot64+3=553\).
Each full three-star occurs in four facets by
Lemma~\ref{lem:multiplicities}. The sum of the facet star counts is
then \(4\cdot553=2212\), exceeding \(3\cdot157+11\cdot158=2209\).

Consequently \(b=0\) on either side. The degree sum \(305\) on a
side of size \(64\) gives the stated counts.
\end{proof}

The degree sequence in Lemma~\ref{lem:degrees} gives
\(\sum_v d_G(v)^2=2(15\cdot4^2+49\cdot5^2)=2930\).
Using \eqref{eq:square} with \(n=7\) and \(e(G)=305\) yields
\begin{equation}\label{eq:globalD}
  D(G)=2.
\end{equation}

\subsection{Opposite facets and the square inequality}

In each \(6\)-facet \(H\) of \(G\), every local degree is in
\(\{3,4,5\}\). For \(i=0,1\) write \(b_{H,i}\) and \(S_{H,i}\)
for the degree-three count and star count on its side inherited from
\(P_i\). Set
\[
  \delta_H=132-e(H),\qquad p_{H,i}=(158-S_{H,i})/3,\qquad
  p_H=p_{H,0}+p_{H,1}.
\]
By \eqref{eq:local-stars}, \(S_{H,i}\equiv2\pmod3\), and
Lemma~\ref{lem:star} gives \(S_{H,i}\le158\). Thus \(p_{H,i}\)
is a nonnegative integer.

\begin{lemma}[Local identities]\label{lem:local-identities}
Every such facet \(H\) satisfies
\begin{equation}\label{eq:local-identities}
  \delta_H\ge0,\qquad
  b_{H,i}=2+2\delta_H-p_{H,i},\qquad D(H)=\delta_H-p_H .
\end{equation}
\end{lemma}

\begin{proof}
Substituting the definitions of \(p_{H,i}\) and \(\delta_H\) into
\eqref{eq:local-stars} gives the identity for \(b_{H,i}\).
For degrees \(3,4,5\), the polynomial \((d-4)(d-5)\) equals
\(2\one_{\{d=3\}}\). Summing and using \eqref{eq:square} in
dimension six gives
\[
  b_{H,0}+b_{H,1}=400-3e(H)+D(H)=4+3\delta_H+D(H).
\]
Comparing with \(b_{H,0}+b_{H,1}=4+4\delta_H-p_H\) yields the
identity for \(D(H)\).

To prove nonnegativity of \(\delta_H\), note that
\(6e(H)-640\le S_{H,i}\le158\), so \(e(H)\le133\).
If \(e(H)=133\), both parity star counts are \(158\), and
\eqref{eq:local-stars} forces \(b_{H,0}=b_{H,1}=0\).
All degrees in \(H\) are then four or five. On restriction to any
\(5\)-facet \(K\) of \(H\), the local degrees belong to \(\{3,4,5\}\),
since deleting one coordinate removes at most one incident edge.
Each of the \(16\) vertices on the chosen side of \(K\) therefore
contributes \(1\pmod3\) to its star count. There are \(40\)
coordinate \(3\)-faces in \(K\), so the number of its faces missing
a star on that side is \(40-S_{K,i}\equiv0\pmod3\).
Star status in a coordinate \(3\)-face depends only on its own edges.
The missing faces in \(K\) are therefore among the two missing faces
on that side of \(H\). It follows that every \(K\) has zero missing
faces, which is impossible since each missing \(3\)-face of \(H\)
lies in three \(5\)-facets. This gives \(e(H)\le132\).
\end{proof}

\begin{lemma}[Opposite facets]\label{lem:opposite}
For opposite \(6\)-facets \(H,H'\),
\begin{equation}\label{eq:opposite}
  3\delta_H+\delta_{H'}\le11+p_H .
\end{equation}
\end{lemma}

\begin{proof}
Let \(L_{H,i}\) count globally degree-four vertices of \(P_i\)
in \(H\), and let \(L_H=L_{H,0}+L_{H,1}\). Let \(t\) count
retained edges between \(P_0\cap V(H)\) and \(P_1\cap V(H')\).
The global degree sums on these two sets are
\[
  e(H)+t=160-L_{H,0},
  \qquad e(H')+t=160-L_{H',1}.
\]
Subtract and use \(L_{H,1}+L_{H',1}=15\) to obtain
\(\delta_H-\delta_{H'}=L_H-15\).
Each local degree-three vertex was globally of degree four, so
\(L_H\ge b_{H,0}+b_{H,1}=4+4\delta_H-p_H\).
Rearranging gives \eqref{eq:opposite}.
\end{proof}

\begin{proposition}[The facet inequality]\label{prop:facet}
Every \(6\)-facet \(H\) of \(G\) satisfies \(2D(H)\le\delta_H\).
\end{proposition}

\begin{proof}
If \(2D(H)>\delta_H\), then \(2p_H<\delta_H\) by
\eqref{eq:local-identities}. The nonnegativity of \(\delta_{H'}\)
and \eqref{eq:opposite} give
\[
  3\delta_H\le11+p_H\le11+\frac{\delta_H-1}{2},
  \qquad 5\delta_H\le21 .
\]
Thus \(1\le\delta_H\le4\), and substituting each possible value
leaves the five rows of Table~\ref{tab:profiles}, up to exchanging
the parity sides. The rows with \(p_H=0\) have star counts
\((158,158)\), and the rows with \(p_H=1\) have star counts
\((158,155)\) up to order. Lemma~\ref{lem:rigid} rules out every row.
\end{proof}

\begin{table}[H]
\centering
\small
\begin{tabular}{@{}rrrccc@{}}
\toprule
\(\delta_H\) & \(p_H\) & \(e(H)\) &
\((b_{H,0},b_{H,1})\) & \((S_{H,0},S_{H,1})\) & Obstruction\\
\midrule
1 & 0 & 131 & \((4,4)\) & \((158,158)\) & \(\mathrm{R0}\)\\
2 & 0 & 130 & \((6,6)\) & \((158,158)\) & \(\mathrm{R0}\)\\
3 & 0 & 129 & \((8,8)\) & \((158,158)\) & \(\mathrm{R0}\)\\
3 & 1 & 129 & \((7,8)\) & \((155,158)\) & \(\mathrm{R1}\)\\
4 & 1 & 128 & \((9,10)\) & \((155,158)\) & \(\mathrm{R1}\)\\
\bottomrule
\end{tabular}
\caption{The five possible violations of the facet inequality, with
the obstruction from Lemma~\ref{lem:rigid}.}\label{tab:profiles}
\end{table}

\subsection{Completion of the proof}

\begin{proof}[Proof of Theorem~\ref{thm:q7}]
Each square lies in five \(6\)-facets and makes the same contribution
to \(D\) in each restriction. Each edge lies in six \(6\)-facets.
Summing Proposition~\ref{prop:facet} over all fourteen facets and using
\eqref{eq:globalD} gives
\[
  20=2\cdot5D(G)=\sum_H 2D(H)
    \le \sum_H\delta_H
    =14\cdot132-6\cdot305=18 ,
\]
a contradiction. The \(304\)-edge construction
\cite{dptv,minamoto} completes the proof.
\end{proof}

\appendix
\section{Certificates for the six-dimensional obstructions}
\label{app:certificates}

\subsection{Symmetry reduction for the two missing stars}
\label{app:orbits}

Represent a \(3\)-face by a pair \((J,z)\), where \(J\) is its
three-coordinate free set, and \(z\) records its fixed bits and has
zeros on \(J\). We use the integers whose set bits encode \(J\) and
\(z\) in Table~\ref{tab:orbits}, numbering coordinates from \(0\).
Fix the first missing face at \((7,0)\) and write \(A=\{0,1,2\}\)
and \(B=\{3,4,5\}\). This normalization is possible by translating
from a vertex of the specified parity side of the first face and
permuting coordinates. Its stabilizer on the chosen even side consists
of independent permutations within \(A\) and \(B\) and the four
even translations supported in \(A\). Its order is \(3!\,3!\,4=144\).

To see that Table~\ref{tab:orbits} covers all second faces, put
\(s=|J\cap A|\). When \(s>0\), an even translation on \(A\) can
change any pattern on \(A\setminus J\), with one free coordinate
absorbing the parity constraint. On \(B\setminus J\) the invariant
is the Hamming weight of the fixed bits. Thus for \(s>0\) and fixed
weight \(w\) on \(B\setminus J\), the orbit has size
\[
  \binom 3s^{\!2}\,2^{3-s}\binom{s}{w}.
\]
The case \(s=3,w=0\) is the first face itself. When \(s=0\),
the fixed pattern on \(A\) has invariant parity, giving two orbits
of size four. The listed orbits are disjoint and cover the other
\(159\) faces.

\begin{table}[htbp]
\centering
\small
\begin{tabular}{@{}rcrrr@{}}
\toprule
Case & Second face \((J,z)\) & Orbit size &
\(\mathrm{R0}\) additions & \(\mathrm{R1}\) additions\\
\midrule
0 & \((7,8)\)   &  3 &    31{,}234 &   115{,}050\\
1 & \((7,24)\)  &  3 &   118{,}954 &   448{,}387\\
2 & \((7,56)\)  &  1 &   929{,}905 &   997{,}572\\
3 & \((11,0)\)  & 18 &   398{,}037 & 1{,}193{,}666\\
4 & \((11,16)\) & 36 &   556{,}394 & 1{,}096{,}140\\
5 & \((11,48)\) & 18 &   755{,}077 & 1{,}160{,}105\\
6 & \((25,0)\)  & 36 &   331{,}182 &   990{,}166\\
7 & \((25,32)\) & 36 &   600{,}382 & 1{,}164{,}830\\
8 & \((56,0)\)  &  4 &   726{,}716 & 1{,}227{,}344\\
9 & \((56,1)\)  &  4 &   391{,}579 &   680{,}982\\
\midrule
 & Total & 159 & 4{,}839{,}460 & 9{,}074{,}242\\
\bottomrule
\end{tabular}
\caption{Orbit representatives and checked clause additions for the
two star-count obstructions.}\label{tab:orbits}
\end{table}

\subsection{Boolean encodings}\label{app:encoding}

For an ambient edge \(e\), let \(x_e\) indicate retention in the graph.
Each coordinate square receives a clause consisting of the negations
of its four edge variables. For a candidate star with edge variables
\(x_1,x_2,x_3\) and indicator \(s\), add the three clauses
\(\neg s\vee x_i\) and the clause
\(s\vee\neg x_1\vee\neg x_2\vee\neg x_3\).
On a given parity side of a \(3\)-face there are four candidate stars.
Their pairwise exclusion clauses enforce at most one, and their
disjunction enforces at least one when required. To bound a vertex
degree between three and five, add the positive clause for each
four-element subset of its incident edges and the negative clause on
all six. For Lemma~\ref{lem:root}, use the degree bounds on the other
even vertices, six unit clauses at the origin, and exactly one even
star in each \(3\)-face. The opposite parity side is unrestricted
except for the square constraints.

For Lemma~\ref{lem:rigid}, impose the degree bounds at every vertex.
At each even vertex define a degree-three indicator \(z\) using
\(\neg z\vee\bigvee_{e\in U}\neg x_e\) for each four-element subset
\(U\) of incident edges and \(z\vee\bigvee_{e\in T}x_e\) for each
three-element subset \(T\). With the degree bounds, \(z\) is true
exactly at degree three. For each orbit representative, forbid the
four even stars on each of the two specified missing faces and
require exactly one even star on all other faces. For each face,
the odd-side missing-star variable \(h\) is encoded by
\(h\vee s_1\vee s_2\vee s_3\vee s_4\) and the four clauses
\(\neg h\vee\neg s_i\), where the \(s_i\) are its odd star
indicators. Count exactly two true \(h\)'s and at least four
true \(z\)'s for \(\mathrm{R0}\), or exactly five true \(h\)'s
and at least eight true \(z\)'s for \(\mathrm{R1}\).

For completeness, consider the auxiliary formula enforcing at most
\(k\) true literals among \(\ell_1,\ldots,\ell_N\), for
\(0<k<N-1\). Use variables \(q_{i,j}\) for \(1\le i\le N-k\)
and \(1\le j\le k\) and the clause families
\[
  \begin{array}{ll}
  \neg\ell_i\vee q_{i,1},&
  \neg q_{i,j}\vee q_{i+1,j}\quad(i<N-k),\\[2pt]
  \neg\ell_{i+j}\vee\neg q_{i,j}\vee q_{i,j+1}\quad(j<k),&
  \neg\ell_{i+k}\vee\neg q_{i,k}.
  \end{array}
\]
Assign \(q_{i,j}\) true if the first \(i+j-1\) inputs contain
at least \(j\) true literals. This extends every assignment satisfying
the bound. Conversely, the first three clause families force the
corresponding variables for all reached prefix thresholds. An induction
on the prefix length proves this claim. A threshold already reached
in the shorter prefix propagates by the second family. A threshold
first reached by the last input follows from the third family or, for
\(j=1\), the first. If the \((k+1)\)-st true input is at position
\(p\), the forced \(q_{p-k,k}\) and \(\ell_p\) contradict the last
family. The bounds \(k=0,N-1,N\) are represented by negative unit
clauses, one all-negative clause, and no clauses, respectively.
Lower bounds use the upper bound on the negated inputs, and exact
counts impose both bounds with separate auxiliary variables.

\subsection{Proof checking and the certificate inventory}
\label{app:replay}

For a proposed derived clause \(C\), the checker assumes the
negation of all its literals and applies unit propagation to the
input and previously verified clauses. A conflict implies \(C\).
Inductively, every verified clause follows from the input. The
checker can use a previously verified clause after deletion, since
the clause remains an input consequence. A verified empty clause
therefore proves unsatisfiability. The checks use unit propagation
with no calls to the solver that generated the traces.

\begin{table}[htbp]
\centering
\small
\begin{tabular}{@{}lrrrr@{}}
\toprule
Obstruction & Inputs & Variables per input & Clauses per input &
Checked additions\\
\midrule
Lemma~\ref{lem:root} & 1 & 832 & 4{,}422 & 291{,}482\\
\(\mathrm{R0}\) & 10 & 2{,}440 & 12{,}974 & 4{,}839{,}460\\
\(\mathrm{R1}\) & 10 & 3{,}438 & 14{,}978 & 9{,}074{,}242\\
\midrule
Total & 21 & & & 14{,}205{,}184\\
\bottomrule
\end{tabular}
\caption{The finite inputs and verified clause additions. The counts
of additions include the empty clause in each refutation.}
\label{tab:checks}
\end{table}

The files for Lemma~\ref{lem:root} have stem
\path{star_q6_def0_root2_flushed}. For orbit \(j\) of
Table~\ref{tab:orbits}, replace \texttt{X} by \(j\) in
\path{star_rigidity_p0_b4_caseX_odd_degrees_cert} for \(\mathrm{R0}\)
or \path{star_rigidity_p1_b8_caseX_odd_degrees_cert} for \(\mathrm{R1}\).
The extensions \texttt{.cnf} and \texttt{.drup} identify the input
and its trace. The archive lists the SHA-256 digests in
\path{star_rigidity_certificate_manifest.json}.
The file \path{local_rigidity_final_audit.json} summarizes the input
reconstruction and proof checks, including the single degree-two
obstruction.

The checkers are
\path{check_drup.ps1} for the single degree-two obstruction and
\path{check_drup_history_fast.ps1} for the twenty orbit inputs.
Run either checker with its \texttt{-Prefix} argument set to the
corresponding input stem. The archived runs checked the additions
listed in Table~\ref{tab:checks}. The matching lower-bound edge list
is available in Minamoto's repository~\cite{minamotodata}.

\subsection*{Data and code availability}
The manuscript, inputs, and verification programs are available in the
\href{https://github.com/zhelizhou/hypercube-c4}{companion repository}.
The \href{https://github.com/zhelizhou/hypercube-c4/tree/main/q7}
{seven-dimensional proof materials} include replay instructions. The
\href{https://github.com/zhelizhou/hypercube-c4/releases/tag/artifacts-v0.1}
{artifact release} contains the \(21\) refutation traces.

\begin{thebibliography}{6}

\bibitem{bhn}
P.~Brass, H.~Harborth, and H.~Nienborg,
\emph{On the maximum number of edges in a \(C_4\)-free subgraph of \(Q_n\)},
J. Graph Theory \textbf{19} (1995), 17--23.
\href{https://doi.org/10.1002/jgt.3190190104}{doi:10.1002/jgt.3190190104}.

\bibitem{dptv}
B.~Derrida, Y.~Pomeau, G.~Toulouse, and J.~Vannimenus,
\emph{Fully frustrated simple cubic lattices and the overblocking effect},
J. Physique \textbf{40} (1979), 617--626.
\href{https://doi.org/10.1051/jphys:01979004006061700}
{doi:10.1051/jphys:01979004006061700}.

\bibitem{erdos}
P.~Erd\H{o}s,
\emph{On some problems in graph theory, combinatorial analysis and combinatorial number theory},
in \emph{Graph Theory and Combinatorics}, B.~Bollob\'as (ed.),
Academic Press, London, 1984, pp.~1--17.

\bibitem{hn}
H.~Harborth and H.~Nienborg,
\emph{Maximum number of edges in a six-cube without four-cycles},
Bull. Inst. Combin. Appl. \textbf{12} (1994), 55--60.

\bibitem{minamoto}
M.~Minamoto,
\emph{\(C_4\)-free subgraphs of the hypercubes \(Q_6\), \(Q_7\), and \(Q_8\):
odd squares, fully frustrated models, and computational structure},
\href{https://arxiv.org/abs/2603.29127v5}{arXiv:2603.29127v5}, 2026.

\bibitem{minamotodata}
M.~Minamoto,
\emph{Supplementary data and code for the hypercube constructions},
\href{https://github.com/minamominamoto/c4free-hypercube}
{GitHub repository}.

\end{thebibliography}

\end{document}
